\chapter{虚拟实验室与评价协议}
\label{ch:sim}

本章把“无实体样机”写成可判定的实验。仿真器内部真值只用于打分。估计器、探针策略和 MPC 只能读取第 \ref{ch:est} 章列出的带噪通道。

\section{双模型原则}

\begin{description}
\item[真值模型 T] 一维或二维多孔介质方程 \eqref{eq:energy}--\eqref{eq:structure}，网格分辨率高于在线模型，参数从先验分布抽样且含空间非均质（局部脂肪或局部厚度 $\pm 10\%$）。
\item[在线模型 E] 两节点或神经算子 \eqref{eq:op}，参数初始化为类别先验，不得读取 T 的参数真值。
\end{description}

若 E 与 T 使用同一套方程和同一套参数，结果只说明离散格式自洽，不说明感知成立。主结果来自 T 与 E 不同的设定。同构试验放在数值检查中，不进入主表。

腔体不必每次都做三维 CFD。协议分两级：

\begin{enumerate}
\item \textbf{边界库。} 用少量 CFD 或实验关联，生成 $h(|v|)$、辐射角系数和蒸汽混合时间常数的查表，作为 T 的边界。
\item \textbf{食品域求解。} 每个控制周期用有限差分或有限元积分食品域。这是计算量的主体。
\end{enumerate}

CFD 的角色是提供可信边界，不是在每个 MPC 迭代里重算腔体。

\section{四套仿真，不混在一个总数里}

仿真按用途分成四套。四套的轨迹都保存，但只有第一套进入主结论。把教师轨迹的条数加进“实验次数”会把数值扩充写成独立证据。

\begin{table}[H]
\centering
\caption{仿真套件与最低规模}
\begin{tabularx}{\textwidth}{lXr}
\toprule
套件 & 回答的问题 & 最低规模 \\
\midrule
S0 格式 & 网格减半后 $\vect{z}$ 的变化是否低于阈值 & 每个原型 6 条 \\
S1 主对照 & B0--B3 在异构真值上的 $e_{\alpha}$ 与信息 & 1\,620 条决策轨迹 \\
S2 边界与失效 & 探针、遮挡、通道缺失、安全边界 & 480 条 \\
S3 教师 & 训练算子并做干预位移测试 & 高保真 360，扩充 $10^4$ 量级 \\
\bottomrule
\end{tabularx}
\end{table}

S1 的格子是：3 个食材原型 $\times$ 3 个厚度 $\times$ 2 个初温 $\times$ 3 个含水偏差 $\times$ 3 个非均质幅度 $\times$ 30 个种子。每格再跑 B0--B3，故 $3\times3\times2\times3\times3\times30=1\,620$ 条决策轨迹，内部含四策略。30 个种子用于分位数，不用单次轨迹代表一格。

S2 不重复 S1 的全因子，只在中等厚度上做：探针库五档、热像遮挡四档、去掉质量或热像、脂肪层有无、开门一次、管温贴近屏障。每项 40 个种子。

S3 的执行器覆盖加热、蒸汽、风速、排汽的拉丁超立方，并强制插入 P1--P4。扩充轨迹由降阶模型生成，日志里标记 \texttt{augmented=true}，不进入 S1 的分母。

\section{主因素}

\begin{table}[H]
\centering
\caption{S1 的因素}
\begin{tabularx}{\textwidth}{llX}
\toprule
因素 & 水平 & 目的 \\
\midrule
食材原型 & 肉糜板、土豆、面团 & 蛋白主导、淀粉主导、高含水多孔 \\
厚度 & 15 mm、30 mm、50 mm & 可辨识性边界 \\
初始温度 & 5 $^\circ$C、20 $^\circ$C & 开环曲线失效的常见原因 \\
初始含水 & 类别中位及 $\pm 8\%$ & 先验误差 \\
非均质 & 均匀、厚度 $\pm 10\%$、脂肪夹层 & 冷点离开几何中心 \\
种子 & 30 & 参数与噪声的独立抽样 \\
\bottomrule
\end{tabularx}
\end{table}

传感器噪声在 S2 中取标定值的 1 倍、2 倍和 4 倍。热像遮挡取 0、15 s、30 s、60 s。这些因素不放进 S1，以免主因子被故障水平稀释。

对照策略固定为四条，不允许只和“不做控制”比较：

\begin{enumerate}
\item \textbf{B0 开环菜单。} 按厚度档选择时间—温度—蒸汽曲线，中途不改。
\item \textbf{B1 腔温 PID + 湿度阈值。} 代表现有蒸烤一体机的被动闭环。
\item \textbf{B2 无探针 MPC。} 与本文相同的 $J$ 和屏障，但 $\cA$ 只有 P0，估计器只做被动更新。
\item \textbf{B3 本文方法。} 估计、P1--P4、因果模型、MPC、屏障。
\end{enumerate}

消融：去掉质量通道、去掉热像、去掉 P4 只留 P1、去掉 $S\to R_{\mathrm{int}}$、去掉屏障（只在仿真中，用于证明屏障会改变越界次数）。

\section{指标}

所有指标在真值模型 T 上计算。

\begin{align}
e_{\alpha}&=\alpha_m^{\mathrm{true}}(t_{\mathrm{stop}})-\alpha_{\mathrm{safe}},\\
e_{T}&=\big\|T^{\mathrm{true}}(\cdot,t)-T^{\mathrm{est}}(\cdot,t)\big\|_{L^2(\Omega_f)},\\
U_T&=\sigma\big(T^{\mathrm{true}}(\cdot,t_{\mathrm{stop}})\big),\\
E&=\int P(t)\,\mathrm{d}t,\\
N_{\mathrm{viol}}&=\text{安全不等式被违反的次数},\\
G_{\mathrm{real}}&=\log\det\vect{\Sigma}(t_-)-\log\det\vect{\Sigma}(t_+).
\end{align}

$t_{\mathrm{stop}}$ 由各策略自己决定。B0 的停止时刻是菜单终点。主要成功准则：

\begin{enumerate}
\item B3 的 $\Prob(e_{\alpha}<0)$ 低于 B0 与 B1；
\item B3 相对 B2 的 $e_T$ 下降，且 $G_{\mathrm{real}}>0$ 的探针比例高于预设值，证明探针真的带来信息；
\item B3 的 $N_{\mathrm{viol}}=0$，同时 B2 在去掉屏障的消融中 $N_{\mathrm{viol}}>0$；
\item 干预测试中，因果模型对 do(蒸汽) 的 $\Delta\alpha_m$ 误差低于黑箱剩余时间模型；
\item 厚度 50 mm 组允许失败，但系统必须输出宽区间并拒绝精确剩余时间。该组用来绘制信息上限，不算总失败。
\end{enumerate}

能量与时间是次要指标。若 B3 用略高能量换来冷点失败率下降，结论按安全优先级书写，不把能耗下降写成必达目标。

\section{统计与图表}

每张主图对应一个已经定义的量：

\begin{itemize}
\item 冷点温度与表面温度的真值/估计分位带；
\item 一次 P4 探针前后 $\vect{\Sigma}$ 的特征值；
\item 路径 A 与路径 B 的偏效应随厚度变化；
\item 四策略的 $e_{\alpha}$ 分布；
\item 信息增益等高线（厚度 $\times$ 含水）；
\item 降级事件时间线。
\end{itemize}

统计用均值、分位数与置信区间。不报告无法由种子复现的“提升约 30\%”。若某次仿真未跑，正文不出现该次的数字。

\section{软件结构}

\begin{center}
\begin{tabularx}{\textwidth}{llX}
\toprule
路径 & 职责 & 依赖 \\
\midrule
\texttt{truth/} & 多孔介质真值步进 & NumPy / FEniCS 或自写差分 \\
\texttt{sensors/} & 加噪声、饱和、遮挡 & 标定参数文件 \\
\texttt{estimate/} & 集合滤波、Fisher & NumPy \\
\texttt{probe/} & 探针库打分 & 估计器接口 \\
\texttt{twin/} & 结构方程与路径分解 & 同一物理核 \\
\texttt{control/} & MPC 与 CBF & CasADi \\
\texttt{operator/} & FNO 或 DeepONet 训练 & PyTorch \\
\texttt{eval/} & 指标与作图 & 只读真值日志 \\
\bottomrule
\end{tabularx}
\end{center}

\texttt{eval/} 与 \texttt{control/} 不共享内存中的真值对象。真值通过日志文件传递，从接口上杜绝偷看。

\section{几何保真度的三档}

三档都要跑，但结论所属的档必须写在图注里。

\begin{table}[H]
\centering
\caption{真值几何}
\begin{tabularx}{\textwidth}{llX}
\toprule
档 & 求解 & 承担的结论 \\
\midrule
T1 & 一维平板或径向，30 个种子的全 S1 & 策略排序与信息上限的主表 \\
T2 & 二维截面，含脂肪夹层与各向异性；每原型每厚度 10 个种子 & 冷点偏移与 T1 排序是否保持 \\
T3 & 腔体边界库：风速三档、蒸汽两档的外部流动，食品域仍用 T1/T2 & $h(v)$ 与角系数的查表，不直接给 $e_{\alpha}$ \\
\bottomrule
\end{tabularx}
\end{table}

在线集合规模 32。离线核对用 128，只比较协方差是否被 32 明显低估。神经算子的训练用 T1 教师，干预测试必须含 T2 的一组，防止算子只记住一维形状。

腔体 CFD 不进入每个控制周期。T3 生成边界库后，T1 与 T2 通过查表读取 $h$、角系数和混合时间。若某档未完成，该档的图位保留坐标，不由低一档的曲线替代。

这些简化限制几何保真度，不改变所检验的信息问题。一维结果不外推为整机三维热场。三维全耦合不在本协议的主结果内。

主图只画已经定义的量：冷点与表面的真值及估计分位带，P4 前后协方差的特征值，路径 A 与路径 B 的偏效应对厚度，四条策略的 $e_{\alpha}$ 分布，厚度—含水平面上的信息增益，以及降级事件的时间线。统计量是均值、分位数和区间。未由种子复现的比例不写入结果。

结果部分只有在 \texttt{eval/} 读完真值日志之后才生成。协议在程序运行前冻结。
